\documentclass[../../../include/open-logic-section]{subfiles}

\begin{document}

\olfileid{sth}{replacement}{intro}
\olsection{سريزه}

تعويض هغه د اصلونو شېما ده چې د $\ZF$ او~$\Z$ تر منځ توپير
راولي۔ له \crefrange{sth:ordinals::chap}{sth:spine::chap} پورې مو ترې
کار واخيست۔ په دې څپرکي کښې به بالاخره دا پوښتنه وڅېړو: آيا تعويض
توجيه کېدلے شي؟

د پوښتنې د کره کولو دپاره بايد پام وکړو چې تعويض په رښتيا هم ډېر
\emph{قوي} دے۔ په دې څپرکي کښې، او بيا په
\olref[sth][card-arithmetic][fix]{sec} کښې، به د دې قوت اندازه
تر يوه حده څرګنده شي۔ خو همدا قوت ښيي چې توجيې ته ئې رښتينې اړتيا شته۔

د توجيې دوه ډولونه به وڅېړو۔ په لنډه توګه، \emph{بهرنۍ} توجيه
هڅه کوي چې يو اصل د هغه د پايلو په مټ توجيه کړي؛ \emph{دننۍ} توجيه
هڅه کوي چې وښيي د بحث رياضيکي مفهومونه پخپله د هغه اصل ملاتړ کوي۔ د
دې توپير معنا به په دې څپرکي کښې لا ښه روښانه شي، خو دا د يوه پراخ
بحث يوازې سر دے۔ د زياتو معلوماتو دپاره په ځانګړي ډول
\citeauthor{Maddy1988a} (\citeyear{Maddy1988a} او
\citeyear{Maddy1988b}) وګورئ۔

\end{document}
